\documentclass[10pt,a4paper]{amsart}
\usepackage[margin=19mm]{geometry}
\usepackage{amsmath,amssymb,mathtools}
\usepackage[scr=rsfs,cal=euler]{mathalfa}
\usepackage{xfrac}
\usepackage[unicode,hidelinks,bookmarks=false]{hyperref}
\newcommand{\ie}{i.e.\ }
\newcommand{\cf}{cf.\ }
\newcommand{\resp}{respectively\ }
\newcommand{\Subsection}{\subsection}
\providecommand{\nobreakdash}{\nobreak\hbox{-}}
\setlength{\emergencystretch}{2em}
\multlinegap0pt
\usepackage{bm}
\usepackage{amssymb}
\usepackage{enumerate}
\usepackage{xcolor}
\usepackage{braket}
\usepackage{algorithm}\usepackage{algpseudocode}
\newtheorem{assumption}{Assumption}
\newtheorem{theorem}{Theorem}
\title{CARE-SAV: A Conditioning-Aware Random-Feature Framework for Energy-Stable Simulation of Gradient Flows}
\author{Bingcheng Hu, Zhaoxiang Li}
\date{Source: arXiv:2608.25300v1 (CC BY 4.0)}
\begin{document}
\maketitle

\noindent\emph{Source and licence.} CARE-SAV: A Conditioning-Aware Random-Feature Framework for Energy-Stable Simulation of Gradient Flows. Version arXiv:2608.25300v1 (CC BY 4.0), distributed by arXiv under the Creative Commons Attribution 4.0 International licence, http://creativecommons.org/licenses/by/4.0/, as recorded in the arXiv licence metadata for this exact version and shown on the abstract page https://arxiv.org/abs/2608.25300v1. Full licence notice: \texttt{licenses/arxiv-2608.25300-CC-BY-4.0.txt}.\par\smallskip

\noindent\textit{Editorial scope.} Counted unit: the article's theorem thm:uus (source0.tex lines 596--603). The article's complete original proof of it is delivered with it (source0.tex lines 604--644, one \texttt{proof} environment of 41 lines). No mathematics is written, repaired or re-derived: the statement and the proof are the article's own lines, in the article's own notation, and no retained line is substituted or re-flowed.  Retained uncounted as the source-local context the proof needs: lines 587--593.  Nothing was omitted from any retained span, so the package carries no omission record.  No retained line contains a citation, so the wrapper carries no bibliography and no bibliography entry.  No retained line contains a figure environment, an image-inclusion command or drawing code, so no image is delivered and the package declares no asset dependency.  Editorial formatting only, in addition: an AMS \texttt{amsart} preamble that restates the article's own declarations and macros used by the retained lines, namely the environments \texttt{assumption}, \texttt{theorem} on separate counters, together with the standard packages those lines need.  The retained environments print in this document as Assumption 1, Theorem 1 in order of appearance, whereas the article prints the same items with the numbering of its own full document.  The complete native source of the article as distributed in the arXiv e-print for this exact version is bundled unchanged as \texttt{sources/208566/source0.tex}.
\begin{assumption}
\label{as1}
  We assume that the reduced operators satisfy
  \begin{equation}
    \boldsymbol{L}_K=\boldsymbol{L}_K^{T}\succeq0,\quad\boldsymbol z^{T}\boldsymbol{G}_K\boldsymbol z\le0 \quad \forall \boldsymbol z\in\mathbb R^K.
  \end{equation}
\end{assumption}

\begin{theorem}[Unconditional Unique Solvability]
\label{thm:uus}
  Let $\boldsymbol b^{n+1/2}=\boldsymbol b_K(\overline{\boldsymbol c}^{\,n+1/2})$ be fixed by explicit extrapolation. Under Assumption~\ref{as1}, the fully discrete CARE-SAV scheme has a unique solution
  $$
    (\boldsymbol c^{n+1},\boldsymbol d^{n+1/2},r^{n+1}) \in\mathbb R^K\times\mathbb R^K\times\mathbb R
  $$
  for every $\Delta t>0$.
\end{theorem}

\begin{proof}
  The fully discrete scheme is equivalent to the block system
  \begin{equation}
    \label{eq:square_system_pf1}
    \begin{bmatrix}
  \boldsymbol{I}_K/\Delta t&-\boldsymbol{G}_K&0\\
  -\boldsymbol{L}_K/2&\boldsymbol{I}_K&-\boldsymbol b^{n+1/2}/2\\
  -(\boldsymbol b^{n+1/2})^{T}/2&0&1
  \end{bmatrix}
  \begin{bmatrix}
  \boldsymbol c^{n+1}\\
  \boldsymbol d^{n+1/2}\\
  r^{n+1}
  \end{bmatrix}
  =
  \begin{bmatrix}
  \boldsymbol c^n/\Delta t\\
  \boldsymbol{L}_K\boldsymbol c^n/2+\boldsymbol b^{n+1/2}r^n/2\\
  r^n-(\boldsymbol b^{n+1/2})^{T}\boldsymbol c^n/2
  \end{bmatrix}.
  \end{equation}
  It is enough to prove that the associated homogeneous system has only the zero solution. Let $(\delta\boldsymbol c,\boldsymbol d,\delta r)$ satisfy
  \begin{equation}
  \label{eq:matrix_pf1}
    \frac{\delta\boldsymbol c}{\Delta t}=\boldsymbol{G}_K\boldsymbol d,
    \quad
    \boldsymbol d=\frac12\boldsymbol{L}_K\delta\boldsymbol c +\frac12\delta r\,\boldsymbol b,
    \quad
    \delta r=\frac12\boldsymbol b^{T}\delta\boldsymbol c,
  \end{equation}
  where $\boldsymbol b=\boldsymbol b^{n+1/2}$. Taking the inner product of the second equation with $\delta\boldsymbol c$ and using the third equation of Eq.~\eqref{eq:matrix_pf1} gives
  \begin{equation}
  \label{eq:uus_pf1}
    \delta\boldsymbol c^{T}\boldsymbol d =\frac12\delta\boldsymbol c^{T}\boldsymbol{L}_K\delta\boldsymbol c +(\delta r)^2 \ge 0.
  \end{equation}
  On the other hand, the first equation of Eq.~\eqref{eq:matrix_pf1} gives
  \begin{equation}
    \delta\boldsymbol c^{T}\boldsymbol d =\Delta t\,\boldsymbol d^{T}\boldsymbol{G}_K\boldsymbol d\le0.
  \end{equation}
  Consequently all nonnegative terms in Eq.~\eqref{eq:uus_pf1} vanish. Since $\boldsymbol{L}_K\succeq0$, we have $\delta\boldsymbol c^{T}\boldsymbol{L}_K\delta\boldsymbol c=0$, so that $\boldsymbol{L}_K\delta\boldsymbol c=0$, and also $\delta r=0$. The second equation in Eq.~\eqref{eq:matrix_pf1} therefore gives $\boldsymbol d=0$, and the first then gives $\delta\boldsymbol c=0$. Thus the homogeneous kernel is trivial, and Eq.~\eqref{eq:square_system_pf1} is nonsingular.
\end{proof}

\end{document}
